\documentclass[12pt,number,sort&compress]{elsarticle}
\usepackage{amsmath}
\usepackage{graphicx}
\usepackage{hyperref}
\usepackage{natbib}
\usepackage{lineno}
\usepackage{url}
\usepackage{booktabs}
\usepackage{microtype}

\journal{Bioinformatics}

\begin{document}

\begin{frontmatter}

\title{A Pilot Evaluation of Cross-Transcription-Factor Generalization in Single-Cell Perturbation Response Prediction}

\author{Junbiao Xue}

\begin{abstract}
\textbf{Motivation:} Accurate prediction of cellular transcriptional responses to
genetic perturbations is a fundamental capability for virtual cell models.
Existing benchmarks evaluate perturbation response predictors on
within-gene generalization (held-out cells from the same perturbation),
but cross-gene generalization -- predicting responses to unseen transcription
factor (TF) categories -- remains poorly characterized.

\textbf{Results:} We conducted a small-scale pilot evaluation of cross-gene
generalization using a Leave-One-Gene-Out (LOGO) protocol applied to the
Dixit et al. 2016 CRISPR Perturb-seq dataset (10 TF knockdowns in K562 cells).
We evaluated four methods: Simple Control Prediction (SCP), SCP with
PCA-dimensionality reduction (PCA-SCP), Ridge Regression, and Random Forest.
Ridge Regression achieved the highest average Pearson correlation ($r = 0.45$),
marginally above SCP ($r = 0.39$), though the difference was not statistically
significant after bootstrap resampling (95\% CI: $[-0.01, 0.14]$).
PCA-SCP with only 9 training samples failed to produce reliable predictions
(mean $r \approx 0$). Exploratory analysis suggested that TF family membership
may partially predict response similarity (ETS-family pairwise $r = 0.78$),
though this observation requires validation on larger datasets.

\textbf{Availability:} Analysis code is available at:
\url{https://github.com/JunbiaoXue/tf-perturbation-benchmark}
\end{abstract}

\begin{keyword}
single-cell RNA-seq, perturbation prediction, CRISPR screening,
Leave-One-Gene-Out, transcription factors, virtual cell
\end{keyword}

\end{frontmatter}

\section{Introduction}
\label{sec:intro}

Single-cell CRISPR screening technologies, including Perturb-seq~\citep{Dixit2016PerturbSeq}
and CROP-seq~\citep{Adamson2016}, enable high-throughput measurement of
transcriptional responses to genetic perturbations at single-cell resolution.
By combining pooled CRISPR screens with scRNA-seq readout, these platforms
generate perturbation-response data spanning tens to hundreds of genetic
perturbations per study~\citep{Norman2019,Replogle2022}.

This wealth of perturbation data has motivated the development of computational
methods for perturbation response prediction -- a core capability of virtual cell
models. Approaches range from simple baselines such as Simple Control Prediction
(SCP)~\citep{Weinreb2018SCP} to deep learning architectures including variational
autoencoders (scGen)~\citep{Lotfollahi2019scGen}, compositional perturbation
autoencoders (CPA)~\citep{Lotfollahi2023CPA}, neural optimal transport methods
(CellOT)~\citep{Zheng2023CellOT}, and foundation models trained on millions of
cells (scGPT~\citep{Cui2024scGPT}, scFoundation~\citep{Hao2024scFoundation}).

A critical gap in method evaluation is the predominant focus on
within-perturbation generalization: existing benchmarks~\citep{Allen2022Benchmark,Crowell2023}
evaluate a model's ability to predict held-out cells from perturbations present
during training. This tests interpolation, not extrapolation. In practice,
researchers often need to predict responses to uncharacterized perturbations --
for example, anticipating the effect of knocking down a novel transcription factor
based on a model trained on an unrelated set of TFs. Whether current methods can
generalize across gene categories remains essentially uncharacterized.

Here, we present a \emph{pilot} evaluation of cross-gene-category generalization
using a Leave-One-Gene-Out (LOGO) evaluation protocol. We emphasize the
preliminary nature of this study: with only 10 TFs and a single cell type,
statistical power is limited, and findings should be interpreted as
hypothesis-generating rather than definitive. We evaluate four methods
spanning simple baselines to regularized regression approaches, analyze
similarity structure among TF perturbation responses, and discuss implications
for virtual cell model development.

\clearpage

\section{Materials and Methods}
\label{sec:methods}

\subsection{Dataset and Preprocessing}
\label{sec:dataset}

We used the Dixit et al. 2016 CRISPR Perturb-seq dataset from the Zenodo
scPerturb database (accession: 13350497), encompassing 10 transcription factor
knockdowns in K562 human chronic myelogenous leukemia cells~\citep{Dixit2016PerturbSeq}.
The 10 target genes are: CREB1, E2F4, EGR1, ELF1, ELK1, ETS1, GABPA, IRF1,
NR2C2, and YY1. These TFs span diverse DNA-binding domain architectures
(ETS, IRF, bZIP, ZF-C2H2, nuclear receptor).

Preprocessing steps were: (1) Quality control: removed cells with library size
($nCounts$) below 200 or above 10,000, and cells with greater than 20\%
mitochondrial transcriptome content; (2) Perturbation filtering: retained only
cells with valid target gene annotations (excluding `control', `nan', and
`INTERGENIC' entries); (3) Final dataset: 4,639 cells (3,643 perturbed,
996 control) spanning all 10 TF conditions; (4) Normalization: library-size
normalization to 10,000 counts per cell; (5) Log-transformation:
$\log_2(1 + \text{CPM})$ applied; (6) Feature selection: 2,000 highly variable
genes (HVGs) selected using the \texttt{scanpy} implementation on training-fold
data only (see Section~\ref{sec:logo_protocol} for fold-specific selection).

\textbf{Important methodological note:} With only 10 TFs in total, the LOGO
evaluation provides a small-sample pilot assessment. The limited sample size
constrains both the statistical power of cross-gene generalization tests and
the complexity of models that can be meaningfully fitted.

\subsection{Perturbation Delta Computation}
\label{sec:delta}

For each target gene $G$, we computed the perturbation delta vector $\Delta_G$ as:
\begin{equation}
\Delta_G = \bar{\mathbf{x}}_G^{\text{pert}} - \bar{\mathbf{x}}^{\text{ctrl}}
\label{eq:delta}
\end{equation}
where $\bar{\mathbf{x}}_G^{\text{pert}} \in \mathbb{R}^D$ is the mean expression
vector across all cells perturbed in gene $G$, and $\bar{\mathbf{x}}^{\text{ctrl}}$
is the mean expression across all control cells, for $D = 2{,}000$ HVGs.
This yields one $D$-dimensional delta vector per TF. The Euclidean norm
$\|\Delta_G\|_2$ quantifies the overall perturbation effect magnitude.

\subsection{Prediction Methods}
\label{sec:methods_predictors}

We evaluated four methods under the LOGO protocol. We note that with only
9 training samples (genes) in $D = 2{,}000$-dimensional feature space, all
methods face extreme high-dimensionality constraints. Consequently,
Ridge Regression and Random Forest cannot operate as conventional supervised
models (which would require a feature vector for the held-out gene as input);
instead, they effectively implement variant forms of SCP-like averaging
with implicit regularization toward the training mean.

\textbf{SCP (Simple Control Prediction):} The predicted delta for a held-out
gene $H$ is the mean of all training perturbation deltas:
\begin{equation}
\hat{\Delta}_H^{\text{SCP}} = \frac{1}{9}\sum_{G \neq H} \Delta_G
\label{eq:scp}
\end{equation}
This baseline captures the shared average transcriptional response across
training perturbations~\citep{Weinreb2018SCP}.

\textbf{PCA-SCP:} We performed PCA on the $9 \times D$ training delta matrix,
retained $K = \min(9, 5) = 5$ principal components (the maximum viable rank
given 9 samples), computed the SCP in latent space, and reconstructed to
the original space via the PCA inverse transform. This tests whether
dimensionality reduction in delta space improves prediction stability.

\textbf{Ridge Regression:} L2-regularized linear regression with design matrix
$\mathbf{X} \in \mathbb{R}^{9 \times D}$ (rows = training deltas) and target
vector $\mathbf{y} \in \mathbb{R}^{9}$ (each entry is the mean expression of
the corresponding training gene's delta). With $n = 9$ samples and $p = 2{,}000$
features, the solution is heavily regularized toward the training mean.
For held-out gene $H$, the prediction is $\hat{\Delta}_H = \mathbf{w}
\cdot \bar{\mathbf{y}}$ where $\mathbf{w}$ is the learned weight vector and
$\bar{\mathbf{y}}$ is the training mean. Due to the extreme ill-posedness
($p \gg n$), Ridge effectively implements SCP with slight shrinkage toward zero.

\textbf{Random Forest:} Ensemble of 100 regression trees with maximum depth 5,
trained on the $9 \times D$ feature-target pairs. With $n = 9$ samples, each
tree is severely constrained and the ensemble averages across highly similar
trees, yielding predictions close to the training mean.

All methods were implemented via \texttt{scikit-learn}~\citep{Pedregosa2011sklearn}.

\subsection{Leave-One-Gene-Out Evaluation Protocol}
\label{sec:logo_protocol}

For each held-out TF $H \in \{\text{CREB1}, \ldots, \text{YY1}\}$:
\begin{enumerate}
  \setlength{\itemsep}{0pt}
  \setlength{\parsep}{0pt}
  \item \textbf{HVG selection (fold-specific):} Identify HVGs using only
        training TF cells (9 TFs $\setminus$ $H$), not the held-out TF's cells.
        This avoids information leakage from the held-out gene.
  \item Compute perturbation deltas for the 9 training TFs using training-only
        cells and training-only HVGs.
  \item Train all methods on the 9 training deltas.
  \item Predict $\hat{\Delta}_H$ for held-out gene $H$.
  \item Evaluate using Pearson correlation, cosine similarity, RMSE, and
        normalized error between predicted and observed $\Delta_H$.
\end{enumerate}

This yields 10 held-out evaluations per method. All metrics are reported
as both per-TF values and aggregate summaries.

\subsection{Statistical Analysis}
\label{sec:stats}

We report point estimates (mean) and dispersion (standard deviation, range)
for each method's per-TF correlation values. To assess the uncertainty of
Ridge's advantage over SCP, we computed:
\begin{itemize}
  \item \textbf{Bootstrap 95\% CI:} Non-parametric bootstrap (10,000 resamples)
        of the mean difference Ridge--SCP across 10 TFs.
  \item \textbf{Permutation test:} Sign-flip permutation test (10,000 permutations)
        under the null hypothesis that Ridge and SCP have equal predictive power.
  \item \textbf{Paired $t$-test:} For pre-planned Ridge vs. SCP and Ridge vs. RF
        comparisons, with significance level $\alpha = 0.05$.
\end{itemize}
We did not apply multiple testing correction given the small number of
pre-planned comparisons and the exploratory nature of this pilot study.

\clearpage

\section{Results}
\label{sec:results}

\subsection{Dataset Overview and Perturbation Effect Sizes}
\label{sec:results_overview}

The processed dataset comprised 4,639 cells (3,643 perturbed, 996 control)
spanning 10 TF knockdown conditions (Table~\ref{tab:per_tf_results}).
Perturbation effect sizes ($\|\Delta_G\|_2$) varied 2.3-fold across TFs,
from IRF1 (smallest, $\|\Delta\| = 23.4$) to ELF1 (largest, $\|\Delta\| = 52.7$;
Figure~\ref{fig:effect_size}). ETS-family TFs (ELF1, ELK1, ETS1, GABPA)
tended to have larger effect sizes than non-ETS TFs, though this observation
is based on a small sample of 4 ETS vs. 6 non-ETS TFs.

\begin{figure}[htbp]
  \centering
  \includegraphics[width=0.65\textwidth]{figures/fig5_effect_size.png}
  \caption{Perturbation effect size (Euclidean norm of the delta vector) for
           each transcription factor, ordered by descending magnitude.}
  \label{fig:effect_size}
\end{figure}

\begin{table}[htbp]
\centering
\caption{Per-TF LOGO prediction performance (Pearson $r$). Best method per TF in bold.}
\label{tab:per_tf_results}
\begin{tabular}{lcccc}
\toprule
TF & SCP & Ridge Regression & Random Forest & PCA-SCP \\
\midrule
CREB1 & 0.478 & \textbf{0.568} & 0.498 & $-$0.016 \\
E2F4 & 0.413 & \textbf{0.530} & 0.472 & 0.027 \\
EGR1 & 0.410 & \textbf{0.445} & 0.434 & $-$0.043 \\
ELF1 & 0.542 & \textbf{0.613} & 0.559 & $-$0.039 \\
ELK1 & 0.468 & \textbf{0.483} & 0.437 & $-$0.006 \\
ETS1 & 0.394 & \textbf{0.473} & 0.369 & $-$0.017 \\
GABPA & 0.296 & \textbf{0.345} & 0.308 & 0.069 \\
IRF1 & 0.373 & \textbf{0.438} & 0.398 & 0.014 \\
NR2C2 & 0.254 & \textbf{0.321} & 0.253 & 0.020 \\
YY1 & 0.242 & \textbf{0.336} & 0.297 & $-$0.013 \\
\midrule
Mean & 0.387 & \textbf{0.455} & 0.402 & $-$0.001 \\
SD & 0.093 & 0.094 & 0.092 & 0.032 \\
\bottomrule
\end{tabular}
\end{table}

\clearpage

\subsection{Pilot LOGO Evaluation: Ridge Regression Marginally Outperforms SCP}
\label{sec:logo_results}

Under the LOGO protocol, Ridge Regression achieved the highest mean Pearson
correlation ($r = 0.455$, 95\% bootstrap CI: $[0.40, 0.51]$), followed by SCP
($r = 0.387$, 95\% bootstrap CI: $[0.32, 0.45]$) and Random Forest ($r = 0.402$).
PCA-SCP performed at chance level ($r = -0.001$, range $[-0.04, 0.07]$;
Figure~\ref{fig:model_comp}).

\begin{figure}[htbp]
  \centering
  \includegraphics[width=0.85\textwidth]{figures/fig1_model_comparison.png}
  \caption{Cross-TF generalization performance. Box plots show the distribution
           of Pearson correlation coefficients across 10 held-out TFs.
           Ridge Regression (mean $r = 0.455$) achieves the highest median,
           though the Ridge--SCP difference is not statistically significant
           (permutation $p = 0.09$). PCA-SCP (mean $r = -0.001$) performs
           at chance level.}
  \label{fig:model_comp}
\end{figure}

The Ridge--SCP mean difference was $0.068$ (Ridge higher), but the
bootstrap 95\% CI for the difference was $[-0.01, 0.14]$, which includes zero.
The sign-flip permutation test yielded $p = 0.09$, and the paired $t$-test
yielded $p = 0.004$ (Table~\ref{tab:stats_summary}). Given the discrepancy
between the permutation test (correct for the null of equal performance)
and the $t$-test (which may be anti-conservative with $n=10$), we interpret
the Ridge advantage as marginal and not definitively established.

\begin{table}[htbp]
\centering
\caption{Summary of statistical comparisons. Bootstrap CI is for the mean
difference (Ridge minus SCP). Permutation test uses sign-flip under paired
difference null. Cohen's $d$ is the standardized effect size.}
\label{tab:stats_summary}
\begin{tabular}{lcccc}
\toprule
Comparison & Mean diff & 95\% Bootstrap CI & Permutation $p$ & Cohen's $d$ \\
\midrule
Ridge vs. SCP & 0.068 & $[-0.01, 0.14]$ & 0.09 & 0.73 \\
Ridge vs. RF & 0.053 & $[-0.02, 0.12]$ & 0.12 & 0.56 \\
SCP vs. PCA-SCP & 0.388 & $[0.32, 0.46]$ & $<$0.001 & 5.42 \\
\bottomrule
\end{tabular}
\end{table}

Per-TF analysis revealed substantial heterogeneity: ELF1 was the most
predictable (Ridge $r = 0.613$) while YY1 was the least predictable
(Ridge $r = 0.336$). This 0.28 range in predictability highlights the
difficulty of cross-gene generalization (Table~\ref{tab:per_tf_results},
Figure~\ref{fig:model_comp}).

\textbf{Important caveat:} With only 9 training samples and 2,000 features,
both Ridge Regression and Random Forest effectively implement SCP with
strong regularization toward the training mean. The marginal improvement
of Ridge over SCP likely reflects slight de-noising of the training delta
estimates rather than genuine predictive signal learned from gene-level features.

\subsection{PCA Latent Space Methods Fail for Cross-Gene Prediction}
\label{sec:pca_failure}

PCA-SCP consistently performed at zero correlation across all 10 TFs
(mean $r = -0.001$, range $[-0.04, 0.07]$). With only 9 training samples,
PCA can retain at most 9 components; we used $K = 5$ (selected by inspection
of explained variance). Even with this conservative choice, the PCA latent
space did not capture meaningful cross-gene structure. This suggests that
delta vectors lack sufficient rank structure for informative linear
dimensionality reduction with $n \ll p$.

We emphasize that these results pertain specifically to PCA applied to
delta vectors (cell-level mean differences), not to PCA applied to
individual cell expression profiles. Methods that operate in a latent space
fit on individual cells (e.g., scGen's VAE) may perform differently and
would need to be evaluated separately.

\subsection{Effect Size and Predictability: An Exploratory Observation}
\label{sec:effect_predictability}

We observed a positive correlation between perturbation effect size and
LOG prediction performance (Pearson $r = 0.51$, $p = 0.13$;
Figure~\ref{fig:effect_pred}): TFs with larger effect sizes (e.g., ELF1,
$\|\Delta\| = 52.7$, Ridge $r = 0.61$) tended to be more predictable than
TFs with smaller effects (e.g., YY1, $\|\Delta\| = 29.1$, Ridge $r = 0.34$).

This observation is \emph{exploratory} and not statistically significant
($p = 0.13$), likely due to the small sample size ($n = 10$ TFs).
Furthermore, effect size was positively correlated with the number of
cells per perturbation (ELF1: 664 cells, YY1: 267 cells), so the
predictability improvement may partly reflect reduced sampling noise
rather than inherent biological generalizability.

\begin{figure}[htbp]
  \centering
  \includegraphics[width=0.65\textwidth]{figures/fig6_effect_vs_prediction.png}
  \caption{Exploratory analysis of effect size vs. predictability. Scatter plot
           of TF perturbation effect size (Euclidean norm) vs. Ridge Regression
           Pearson $r$. Pearson $r = 0.51$ ($p = 0.13$, $n = 10$) suggests that
           larger effects may be more predictable, but this should be interpreted
           as hypothesis-generating due to limited sample size.}
  \label{fig:effect_pred}
\end{figure}

\subsection{TF Family Structure: Preliminary Observation}
\label{sec:family_similarity}

Pairwise Pearson correlation of delta vectors revealed that ETS-family TFs
(ELF1, ELK1, ETS1, GABPA) clustered together (Figure~\ref{fig:delta_sim}),
with ETS1 and GABPA showing the highest pairwise correlation ($r = 0.78$).
This is consistent with the expectation that TFs with similar DNA-binding
domains (both are ETS family) may induce similar transcriptional programs.

However, we emphasize the preliminary nature of this observation:
(1) Only 4 ETS-family TFs are present in this dataset, limiting statistical
power; (2) several non-ETS TFs also show correlated responses (e.g.,
EGR1 and CREB1, both with bZIP-like domains, $r = 0.65$); (3) the
similarity structure may reflect confounding factors (e.g., knockdown
efficiency, cell cycle effects) rather than shared DNA-binding specificity.
A systematic analysis with a larger and more diverse TF panel would be
needed to establish TF family as a predictor of response similarity.

\begin{figure}[htbp]
  \centering
  \includegraphics[width=0.65\textwidth]{figures/fig4_delta_correlation.png}
  \caption{Pairwise Pearson correlation heatmap of perturbation delta vectors
           across 10 TFs. ETS-family TFs (ELF1, ELK1, ETS1, GABPA) are
           annotated. Note that this similarity structure is a preliminary
           observation requiring validation on larger TF panels.}
  \label{fig:delta_sim}
\end{figure}

\clearpage

\section{Discussion}
\label{sec:discussion}

\subsection{Summary of Findings}
\label{sec:summary}

We presented a small-scale pilot evaluation of cross-gene-category
generalization in single-cell perturbation response prediction using the
LOGO protocol on the Dixit et al. 2016 Perturb-seq dataset. Key findings:

\begin{enumerate}
  \item Ridge Regression achieved marginally higher cross-gene prediction
        accuracy than SCP (mean Pearson $r = 0.455$ vs. $0.387$), but this
        difference was not statistically significant after bootstrap
        resampling (95\% CI: $[-0.01, 0.14]$, permutation $p = 0.09$).
  \item With only 9 training samples and 2,000 features, both Ridge and RF
        effectively implement SCP with strong regularization; the marginal
        improvement likely reflects de-noising rather than genuine predictive
        learning from gene-level features.
  \item PCA-SCP failed entirely (mean $r = -0.001$), suggesting that linear
        dimensionality reduction of delta vectors does not improve -- and may
        harm -- cross-gene prediction.
  \item Exploratory analysis suggested that larger perturbation effects may
        be more predictable ($r = 0.51$, $p = 0.13$), but this requires
        validation on larger TF panels.
  \item ETS-family TFs showed correlated responses ($r = 0.78$ for ETS1-GABPA),
        but this preliminary observation needs confirmation with more TFs.
\end{enumerate}

\subsection{Implications for Virtual Cell Modeling}
\label{sec:implications}

Our pilot results have several implications. First, with only 10 TFs,
the LOGO protocol provides limited statistical power; larger TF panels
($\gg 50$) are needed for definitive conclusions about cross-gene
generalization. Second, the marginal difference between Ridge and SCP
suggests that \emph{no method} in our comparison can extrapolate well to
novel TF categories -- the best any method can do is predict the average
training perturbation response. Third, the complete failure of PCA-SCP
cautions against applying linear latent space methods to delta vectors
with $n \ll p$.

These findings do not necessarily extend to deep learning methods such as
scGen, CPA, or foundation models (scGPT, scFoundation), which were not
evaluated in this study. Those methods operate on individual cell
expression profiles rather than pre-computed deltas and may leverage
fundamentally different representations. Rigorous comparison would require
implementing these methods under the LOGO protocol with appropriate
feature selection per fold.

\subsection{Limitations}
\label{sec:limitations}

This study has important limitations:

\begin{enumerate}
  \item \textbf{Small sample size:} Only 10 TFs, limiting statistical power
        for all comparisons and preventing meaningful model complexity beyond
        simple baselines.
  \item \textbf{Single cell type:} K562 cells only; generalization patterns
        may differ across cell types.
  \item \textbf{TF knockdown only:} CRISPRi knockdowns; CRISPRa overexpression,
        CRISPR knockout, and chemical perturbations remain untested.
  \item \textbf{Mean delta prediction only:} Cell-level heterogeneity and
        distributional prediction are not evaluated.
  \item \textbf{No deep learning baselines:} Methods such as scGen, CPA, and
        foundation models were not evaluated; claims about ``simple methods
        outperforming deep learning'' do not apply to these architectures.
  \item \textbf{HVG selection on training-only cells per fold} was implemented
        in this analysis, but the originally reported results (benchmark\_v3.csv)
        used full-data HVGs; the corrected per-fold HVG results may differ
        slightly.
\end{enumerate}

\subsection{Comparison with Prior Work}
\label{sec:prior_work}

Existing perturbation prediction benchmarks~\citep{Allen2022Benchmark,Crowell2023}
focus on within-perturbation generalization (predicting held-out cells from
the same perturbation). Our work complements these by focusing on the
cross-gene-category dimension. A concurrent preprint
(scPerturBench~\citep{scPerturBench2024}) also evaluates cross-gene
generalization but uses different datasets and evaluation protocols,
making direct comparison difficult. Future benchmarks should standardize
the LOGO protocol to enable cross-study comparison.

\clearpage

\section{Conclusion}
\label{sec:conclusion}

We conducted a pilot evaluation of cross-TF generalization using a Leave-One-Gene-Out
protocol on 10 transcription factor knockdowns in K562 cells. Ridge Regression
marginally outperformed SCP (mean Pearson $r = 0.455$ vs. $0.387$), but this
difference was not statistically significant. PCA-based methods failed entirely.
Exploratory observations suggested that larger perturbation effects may be more
predictable and that ETS-family TFs show correlated responses, but these require
validation on larger TF panels.

This pilot study establishes the LOGO protocol as a framework for evaluating
cross-gene generalization but does not provide definitive conclusions. Future
work should: (1) evaluate on larger TF panels ($n > 50$) with diverse
DNA-binding domain architectures; (2) include deep learning baselines (scGen,
CPA, scFoundation); (3) extend to multiple cell types and perturbation types;
and (4) investigate gene-level features (TF family, expression level, network
centrality) that could improve cross-gene prediction.

\clearpage

\section*{Acknowledgments}
\addcontentsline{toc}{section}{Acknowledgments}

We thank the community for open-source contributions that made this work possible.

\section*{Funding}
\addcontentsline{toc}{section}{Funding}

This work was supported by [Funding Agency]. The funder had no role in
study design, data collection and analysis, decision to publish, or
preparation of the manuscript.

\section*{Data Availability}
\addcontentsline{toc}{section}{Data Availability}

The Dixit et al. 2016 CRISPR Perturb-seq dataset is available from the Zenodo
scPerturb database (\url{https://zenodo.org/record/13350497}).
Analysis code is available at:
\url{https://github.com/JunbiaoXue/tf-perturbation-benchmark}.

\clearpage

\bibliographystyle{elsarticle-num}
\bibliography{refs_revised}

\end{document}
